\begin{resumo}[Abstract]
\begin{otherlanguage*}{english}
This document presents the formal specification of the Database Design and Data Modeling for the Intelligent Adverse Event Management System based on the Global Trigger Tool methodology (SIGEA-GTT), developed at the Federal University of Sergipe (UFS). The system operates across a tripartite data frontier comprising strict institutional identity access control, standardized clinical adverse event taxonomy (IHI and NCC MERP), and pedagogical management of academic cohorts, evaluations, and simulated medical records. The data architecture adopts a hybrid relational-document model implemented in PostgreSQL 16 LTS DBMS, utilizing strictly normalized tables (3NF/BCNF) with universal UUID v4 keys for relational integrity, combined with semi-structured \texttt{JSONB} attributes to represent high-fidelity synthetic electronic health records and quality management tool artifacts (Ishikawa, 5W2H, GUT, and PDCA). The Conceptual Entity-Relationship Model based on Peter Chen's notation, the Logical Relational Diagram in Crow's Foot notation, the exhaustive data dictionary for all eight physical entities, the formal JSON Schemas, B-Tree and GIN indexing strategies, domain integrity constraints, and the bidirectional traceability matrix mapping the system requirements are presented, ensuring ethical isolation and strict compliance with the Brazilian General Data Protection Law (LGPD).

\vspace{\onelineskip}
\noindent
\textbf{Keywords}: Relational Database. Entity-Relationship Model. PostgreSQL. JSONB. Global Trigger Tool. Software Engineering.
\end{otherlanguage*}
\end{resumo}
